\documentclass[10pt,letterpaper]{article}
\usepackage{cornell-notes}

\title{Chapter 7: Virtualization and the Cloud}
\author{}
\date{}

\setDocTitle{Chapter 7: Virtualization and the Cloud}
\setDocSubtitle{Operating Systems}
\setDocVersion{v1.0}
\setDocStatus{Study Companion}
\setDocOwner{Operating Systems Cornell Notes}
\setDocAuthor{}
\setDocDate{\today}
\setDocSummary{Cornell-format study and review notes for Chapter 7: Virtualization and the Cloud.}

\setCornellCollection{Operating Systems Cornell Notes}
\setCornellUnitType{Chapter}
\setCornellUnitNumber{7}
\setCornellUnitTitle{Virtualization and the Cloud}
\setCornellCompanionLabel{Study and review companion}

\hypersetup{pdftitle={Chapter 7: Virtualization and the Cloud},pdfsubject={Cornell notes},pdfauthor={}}

\begin{document}
\maketitle
\section*{Learning Objectives}
\begin{studybox}{By the end of this chapter, you should be able to}
\begin{enumerate}
  \item Explain the historical motivations for virtual machines and their resurgence.
  \item Apply the sensitive-instruction and trap-and-emulate requirements for classical virtualization.
  \item Compare type 1 and type 2 hypervisors.
  \item Explain binary translation, paravirtualization, hardware assistance, and their costs.
  \item Trace shadow or nested address translation and representative I/O virtualization paths.
  \item Distinguish virtual appliances, multicore guests, licensing concerns, and cloud service models.
  \item Explain migration and checkpointing state transitions.
  \item Relate the VMware designs to the general virtualization mechanisms.
\end{enumerate}
\end{studybox}

\begin{studybox}{Section Map}
\hyperlink{sec-7-1}{7.1} \quad \hyperlink{sec-7-2}{7.2} \quad \hyperlink{sec-7-3}{7.3} \quad \hyperlink{sec-7-4}{7.4} \quad \hyperlink{sec-7-5}{7.5} \quad \hyperlink{sec-7-6}{7.6} \quad \hyperlink{sec-7-7}{7.7} \quad \hyperlink{sec-7-8}{7.8} \quad \hyperlink{sec-7-9}{7.9} \quad \hyperlink{sec-7-10}{7.10} \quad \hyperlink{sec-7-11}{7.11} \quad \hyperlink{sec-7-12}{7.12} \quad \hyperlink{sec-7-13}{7.13}
\end{studybox}

\section*{Cornell Cue-and-Notes Area}
\small
\begin{longtable}{C N}
\rowcolor{navy}
\color{white}\textbf{Cue / Recall Prompt} &
\color{white}\textbf{Notes}\\
\endfirsthead
\rowcolor{navy}
\color{white}\textbf{Cue / Recall Prompt} &
\color{white}\textbf{Notes (continued)}\\
\endhead
\rowcolor{ice}\multicolumn{2}{l}{\hypertarget{sec-7-1}{}\textbf{Section 7.1: History}}\\*\hline
\textbf{Why virtualization returned} & \textcolor{legacy}{\textbf{Historical or legacy context:}} Early virtual machines partitioned costly mainframes. Later commodity processors, server consolidation, isolation, testing, compatibility, and data-center management renewed demand.\\\hline
\rowcolor{ice}\multicolumn{2}{l}{\hypertarget{sec-7-2}{}\textbf{Section 7.2: Requirements for Virtualization}}\\*\hline
\textbf{Trap and emulate} & A virtual-machine monitor must retain control, present an environment equivalent to hardware, and run most guest instructions directly. Sensitive operations must trap or otherwise be intercepted.\\\hline
\rowcolor{ice}\multicolumn{2}{l}{\hypertarget{sec-7-3}{}\textbf{Section 7.3: Type 1 and Type 2 Hypervisors}}\\*\hline
\textbf{Hypervisor placement} & A type 1 hypervisor runs directly on hardware and manages guests. A type 2 monitor runs as or within a host operating-system process and relies on host drivers and services.\\\hline
\rowcolor{ice}\multicolumn{2}{l}{\hypertarget{sec-7-4}{}\textbf{Section 7.4: Techniques for Efficient Virtualization}}\\*\hline
\textbf{7.4.1 Unvirtualizable instructions} & When sensitive instructions do not trap, a monitor can scan and translate basic blocks, replace operations with safe sequences, or modify a guest to issue explicit hypercalls. Hardware extensions later add controlled guest execution modes.\\\hline
\textbf{7.4.2 Cost} & Direct execution is cheap; exits, world switches, translation, emulation, interrupt delivery, and address translation add overhead. Optimization focuses on reducing the frequency and cost of these transitions.\\\hline
\rowcolor{ice}\multicolumn{2}{l}{\hypertarget{sec-7-5}{}\textbf{Section 7.5: Are Hypervisors Microkernels Done Right?}}\\*\hline
\textbf{Comparison} & Both hypervisors and microkernels minimize privileged mechanisms and isolate components. A hypervisor exports a machine interface to unmodified operating systems, whereas a microkernel exports IPC and kernel abstractions to cooperating servers.\\\hline
\rowcolor{ice}\multicolumn{2}{l}{\hypertarget{sec-7-6}{}\textbf{Section 7.6: Memory Virtualization}}\\*\hline
\textbf{Two translation levels} & A guest maps virtual to guest-physical addresses, while the hypervisor maps guest-physical to machine frames. Shadow tables combine mappings in software; nested paging lets hardware walk both levels.\\\hline
\textbf{Memory overcommit} & Balloon drivers, page sharing, reclamation, and hypervisor swapping can support aggregate guest allocations larger than RAM, but reclamation may interact badly with guest replacement decisions.\\\hline
\rowcolor{ice}\multicolumn{2}{l}{\hypertarget{sec-7-7}{}\textbf{Section 7.7: I/O Virtualization}}\\*\hline
\textbf{Device paths} & A monitor may emulate familiar hardware, use paravirtual front-end and back-end drivers, or assign devices with hardware isolation. Compatibility, performance, sharing, and migration pull in different directions.\\\hline
\rowcolor{ice}\multicolumn{2}{l}{\hypertarget{sec-7-8}{}\textbf{Section 7.8: Virtual Appliances}}\\*\hline
\textbf{Packaged machine images} & A virtual appliance bundles an application with a configured guest environment. It simplifies deployment and rollback but creates image maintenance, patching, provenance, and data-separation responsibilities.\\\hline
\rowcolor{ice}\multicolumn{2}{l}{\hypertarget{sec-7-9}{}\textbf{Section 7.9: Virtual Machines on Multicore CPUs}}\\*\hline
\textbf{Virtual CPU scheduling} & The hypervisor maps virtual CPUs to physical CPUs, preserves time and interrupt semantics, and may co-schedule related virtual CPUs. Oversubscription and topology awareness affect latency.\\\hline
\rowcolor{ice}\multicolumn{2}{l}{\hypertarget{sec-7-10}{}\textbf{Section 7.10: Licensing Issues}}\\*\hline
\textbf{Machine identity and mobility} & \textcolor{legacy}{\textbf{Historical or legacy context:}} Software licenses tied to physical processors, machines, or installations become ambiguous when images are copied, paused, cloned, or migrated.\\\hline
\rowcolor{ice}\multicolumn{2}{l}{\hypertarget{sec-7-11}{}\textbf{Section 7.11: Clouds}}\\*\hline
\textbf{Cloud abstraction} & Cloud platforms pool infrastructure and expose elastic, measured resources through network APIs. Virtualization is a common enabling mechanism but not the definition of cloud computing.\\\hline
\textbf{7.11.1 Service models} & Infrastructure service exposes virtual machines and storage; platform service exposes managed runtimes; software service exposes complete applications. Each level shifts control and operational responsibility.\\\hline
\textbf{7.11.2 Migration} & Pre-copy migration transfers memory while the VM runs, then briefly stops it to send remaining dirty state and resume at the destination. Shared storage or storage migration must preserve device state and connectivity.\\\hline
\textbf{7.11.3 Checkpointing} & A checkpoint captures processor, memory, and device state so execution can restart later. Consistent snapshots require coordination with in-flight I/O, external communication, and persistent storage.\\\hline
\rowcolor{ice}\multicolumn{2}{l}{\hypertarget{sec-7-12}{}\textbf{Section 7.12: Case Study: VMware}}\\*\hline
\textbf{7.12.1 Early history} & \textcolor{legacy}{\textbf{Historical or legacy context:}} The design targeted unmodified commodity operating systems on x86 hardware that was not classically virtualizable, motivating a hosted monitor and dynamic translation.\\\hline
\textbf{7.12.2 Workstation} & \textcolor{legacy}{\textbf{Historical or legacy context:}} VMware Workstation used a host operating system for device support and user experience while a virtual-machine monitor isolated and executed each guest.\\\hline
\textbf{7.12.3 x86 challenges} & \textcolor{legacy}{\textbf{Historical or legacy context:}} Some privilege-sensitive x86 instructions behaved differently without trapping, and address-space control had to separate host, monitor, and guest contexts.\\\hline
\textbf{7.12.4 Solution overview} & Dynamic binary translation rewrote sensitive kernel instruction sequences; direct execution handled safe user code; world switches moved between host and monitor contexts; shadow page tables virtualized memory.\\\hline
\textbf{7.12.5 Evolution} & \textcolor{legacy}{\textbf{Historical or legacy context:}} Later processors introduced virtualization modes and nested page-table support, reducing dependence on translation and software-maintained shadow mappings.\\\hline
\textbf{7.12.6 ESX Server} & \textcolor{legacy}{\textbf{Historical or legacy context:}} ESX moved to a type 1 design optimized for virtual machines, with dedicated scheduling, memory overcommit, I/O paths, VMFS storage, and coordinated live migration capabilities.\\\hline
\rowcolor{ice}\multicolumn{2}{l}{\hypertarget{sec-7-13}{}\textbf{Section 7.13: Research on Virtualization and the Cloud}}\\*\hline
\textbf{Research themes} & Research spans isolation, nested virtualization, migration, deduplication, scheduling, fault resilience, debugging, cloud workload management, security, and energy reduction.\\\hline
\end{longtable}
\normalsize

\section*{Chapter Synthesis}
\begin{studybox}{Chapter Summary}
Virtualization inserts a control layer that safely multiplexes hardware while presenting complete machine interfaces to guests. Efficient designs combine direct execution with controlled interception, two-level memory translation, and emulated, paravirtual, or assigned I/O. The same mechanisms enable packaging, consolidation, migration, checkpointing, and many cloud infrastructure services.
\end{studybox}

\begin{studybox}{Key Relationships}
\begin{itemize}
  \item CPU virtualization protects privileged state; memory virtualization protects mappings; I/O virtualization protects and multiplexes devices.
  \item A VM exit is analogous to a trap, but it transfers control from a guest context to the hypervisor.
  \item Migration depends on checkpoint-quality capture plus a transport and destination resource plan.
  \item Hardware assistance changes mechanism cost without removing the need for hypervisor scheduling and isolation policy.
\end{itemize}
\end{studybox}

\begin{studybox}{Design Tradeoffs}
\begin{itemize}
  \item Emulated devices maximize compatibility, while paravirtual or assigned devices improve performance.
  \item A hosted hypervisor reuses mature drivers; a bare-metal hypervisor gains control and lower I/O overhead.
  \item Memory overcommit raises utilization but can amplify paging and unpredictable latency.
  \item Pre-copy migration reduces downtime but may retransmit pages dirtied repeatedly.
\end{itemize}
\end{studybox}

\begin{studybox}{Common Confusions}
\begin{itemize}
  \item A hypervisor virtualizes a machine interface; a container shares the host kernel and isolates at a higher level.
  \item Cloud computing can use virtualization, but pooling, elasticity, APIs, and metering are separate cloud properties.
  \item Guest-physical addresses are an intermediate abstraction and need not equal machine addresses.
\end{itemize}
\end{studybox}

\section*{Key Terms}
\small
\begin{longtable}{C N}
\rowcolor{navy}\color{white}\textbf{Term} & \color{white}\textbf{Concise definition}\\
\endfirsthead
\rowcolor{navy}\color{white}\textbf{Term} & \color{white}\textbf{Concise definition (continued)}\\
\endhead
\textbf{Hypervisor} & Control layer that creates and manages virtual machines.\\\hline
\textbf{Type 1 hypervisor} & Hypervisor executing directly on hardware.\\\hline
\textbf{Type 2 hypervisor} & Hosted monitor relying on a host operating system.\\\hline
\textbf{Trap and emulate} & Execute safe instructions directly and intercept sensitive operations.\\\hline
\textbf{Binary translation} & Runtime rewriting of guest instruction sequences into safe equivalents.\\\hline
\textbf{Paravirtualization} & Guest modification to invoke hypervisor operations explicitly.\\\hline
\textbf{VM exit} & Transfer from guest execution to hypervisor control.\\\hline
\textbf{Shadow page table} & Hypervisor-maintained combined guest-to-machine mapping.\\\hline
\textbf{Nested paging} & Hardware traversal of guest and hypervisor translation structures.\\\hline
\textbf{Ballooning} & Guest-cooperative memory reclamation initiated by the hypervisor.\\\hline
\textbf{Virtual appliance} & Preconfigured application and guest environment packaged as an image.\\\hline
\textbf{Live migration} & Movement of a running VM with bounded service interruption.\\\hline
\textbf{Checkpoint} & Captured execution state from which a VM can resume.\\\hline
\textbf{IaaS} & Cloud service exposing infrastructure resources such as virtual machines.\\\hline
\end{longtable}
\normalsize

\enlargethispage{2\baselineskip}
\begin{studybox}{Active-Recall Review}
\small
\begin{enumerate}
  \item What properties must a classical virtual-machine monitor provide?
  \item How do type 1 and type 2 hypervisors differ?
  \item Why did early x86 virtualization need binary translation?
  \item How do shadow and nested page tables differ?
  \item What are the major I/O virtualization choices?
  \item Why can memory overcommit trigger double paging?
  \item Which state must migration preserve besides memory contents?
  \item How does checkpointing differ from an ordinary memory copy?
  \item Which responsibilities shift among IaaS, PaaS, and SaaS?
\end{enumerate}
\end{studybox}

\par\medskip
{\color{navy}\bfseries Next-Chapter Connections}\par\smallskip
Chapter 8 expands resource coordination from virtual machines to shared-memory multiprocessors, message-passing multicomputers, and distributed systems.
\normalsize
\end{document}
